\documentclass[a4paper,10pt]{article}
\usepackage[utf8]{inputenc}
\usepackage[T1]{fontenc}
\usepackage{lmodern}
\usepackage[margin=1.4cm, top=1.2cm, bottom=1.2cm]{geometry}
\usepackage{hyperref}
\usepackage{xcolor}
\usepackage{enumitem}
\usepackage{graphicx}
\usepackage{titlesec}

% Minimal clean styling
\hypersetup{colorlinks=true, urlcolor=blue!70!black, linkcolor=blue!70!black}
\pagestyle{empty}
\setlength{\parindent}{0pt}

% Section formatting
\titleformat{\section}{\large\bfseries\color{black}}{}{0em}{}[\titlerule]
\titlespacing*{\section}{0pt}{6pt}{4pt}

% Colors
\definecolor{accent}{RGB}{20, 50, 120}

% Role command
\newcommand{\role}[4]{%
  \vspace{2pt}\noindent
  \begin{minipage}[t]{0.70\textwidth}\raggedright\textbf{#1} \textbullet\ \textit{#3}\end{minipage}%
  \hfill
  \begin{minipage}[t]{0.28\textwidth}\raggedleft\footnotesize\color{gray}#2 \textbullet\ #4\end{minipage}%
  \vspace{1pt}%
}
\setlist[itemize]{leftmargin=1.2em, itemsep=1.5pt, topsep=2pt, parsep=0pt}

\begin{document}

% ================= HEADER =================
\begin{minipage}[c]{0.80\textwidth}
  {\Huge\bfseries Prabhavit Gupta}\\[3pt]
  {\color{accent}\large AI Engineer $\vert$ Agentic AI, LLM Orchestration, MCP Integrations, RAG $\&$ RAGAS}\\[6pt]
  {\footnotesize +39 351 610 7807 ~$\vert$~
  \href{mailto:prabhavitg@gmail.com}{prabhavitg@gmail.com} ~$\vert$~
  \href{https://www.linkedin.com/in/prabhavit/}{linkedin.com/in/prabhavit} ~$\vert$~
  \href{https://github.com/prab-gupta}{github.com/prab-gupta} ~$\vert$~
  Venezia / Milano, Italia (Residente Permanente UE $\vert$ Full Remote $\vert$ Disponibilità Immediata)}
\end{minipage}\hfill

\vspace{6pt}

% ================= EXECUTIVE SUMMARY =================
\section{Profilo Professionale / Career Summary}
\begin{itemize}
  \item \textbf{AI Engineer specializzato in Sistemi Agentici ed Architetture LLM End-to-End:} Progettazione ed implementazione di pipeline basate su Large Language Models (OpenAI, Anthropic, vLLM, Ollama), orchestrazione multi-agente (LangGraph, LangChain, LlamaIndex), intent classification, routing semantico e query rewriting/expansion in produzione.
  \item \textbf{Integrazione Avanzata Tool Calling, MCP e Sistemi RAG Ibridi:} Esperienza diretta nell'adozione del \textit{Model Context Protocol (MCP)} per collegare agenti autonomi a database, API aziendali e servizi esterni; sviluppo di architetture RAG avanzate con hybrid search (BM25 + dense embedding), reranking semantico, Knowledge Graph (GraphRAG / Neo4j) e chunking dinamico per dati non strutturati.
  \item \textbf{Evaluation Frameworks, Sicurezza e Validazione RAGAS:} Implementazione sistematica di metriche di valutazione quantitativa (RAGAS: context precision/recall, answer relevancy, faithfulness) per eliminare allucinazioni; implementazione di NeMo Guardrails, regression test per prompt e pipeline AI scalabili ed orientate alla produzione business.
\end{itemize}

% ================= TECHNICAL SKILLS =================
\section{Competenze Tecniche / Technical Skills}
\begin{itemize}[leftmargin=1.2em, itemsep=2pt]
  \item \textbf{LLM \& Agent Orchestration:} Agentic AI Design Patterns, LangGraph, LangChain, LlamaIndex, Multi-Agent Systems, Autonomous Tool Calling, ReAct, Intent Classification, Query Routing \& Expansion.
  \item \textbf{Protocolli \& Integrazione Tool:} Model Context Protocol (MCP - Anthropic/Open Source), Custom MCP Server/Client Development, Function Calling, REST APIs, JSON Schema Enforcement, Structured Outputs (Pydantic).
  \item \textbf{RAG, Retrieval \& Knowledge Graphs:} Retrieval-Augmented Generation (RAG), Hybrid Search (Dense + Sparse/BM25), Cross-Encoder Reranking, Vector DBs (Qdrant, Pinecone, Chroma, pgvector), GraphRAG, Neo4j, Entity \& Relationship Extraction.
  \item \textbf{AI Models \& Deployment:} OpenAI (GPT-4o), Anthropic Claude 3.5 Sonnet, open-source LLMs via vLLM, Ollama, Hugging Face Transformers, Token Optimization, Cost/Latency Tradeoff Analysis.
  \item \textbf{Evaluation \& Sicurezza LLM:} RAGAS Framework, DeepEval, Prompt Regression Testing, Hallucination Detection, NeMo Guardrails, Input/Output Filtering, PII Masking, OWASP Top 10 for LLMs.
  \item \textbf{Linguaggi \& Backend Engineering:} Python (FastAPI, AsyncIO, PyDantic, PyTest), TypeScript / Node.js, SQL (PostgreSQL), Docker, CI/CD, Git, Linux.
\end{itemize}

% ================= PROFESSIONAL EXPERIENCE =================
\section{Esperienza Professionale / Work Experience}

\role{AI Engineer \& Solutions Architect}{2024 -- Presente}{Indipendente / Consulenza AI Enterprise}{Venezia \& Milano, IT}
\begin{itemize}
  \item Progettazione e rilascio di architetture agentiche end-to-end con orchestrazione complessa via LangGraph e protocollo MCP (Model Context Protocol), connettendo agenti decisionali ad ERP, database vettoriali ed API di terze parti.
  \item Costruzione di pipeline RAG ibride ad alte prestazioni su corpus documentali eterogenei: implementazione di chunking semantico, generazione embedding bilingue (italiano/inglese), hybrid retrieval (BM25 + vettoriale) e Cohere reranker, riducendo l'hallucination rate di oltre il 40\%.
  \item Integrazione di Knowledge Graph ed approcci GraphRAG per mappare entità e relazioni complesse da documentazione non strutturata, migliorando l'accuratezza di retrieval per query contestuali multi-step.
  \item Implementazione di benchmark quantitativi continui con RAGAS (faithfulness, answer relevance, context precision), automatizzando regression test per prompt prima di ogni rilascio in esercizio.
  \item Configurazione di guardrails di sicurezza per il filtraggio di prompt injection, mascheramento PII e validazione dello schema JSON di output tramite Pydantic in scenari enterprise.
\end{itemize}

\vspace{3pt}
\role{Full-Stack AI Software Engineer}{2023 -- 2024}{Tech Automation Systems}{Milano / Remoto, IT}
\begin{itemize}
  \item Sviluppo di microservizi ad alto throughput in Python (FastAPI) e TypeScript dedicati all'elaborazione automatizzata di flussi informativi aziendali ed integrazione LLM.
  \item Ingegnerizzazione di pipeline di data ingestion per documenti non strutturati (PDF, contratti, schede tecniche) con parsing, normalizzazione, metadati arricchiti ed indicizzazione su vector store (Qdrant e pgvector).
  \item Ottimizzazione della latenza e dei costi inferenziali combinando modelli proprietari (OpenAI) con runtime locali open source (Ollama, vLLM) per task specifici di classificazione ed estrazione entità.
  \item Collaborazione diretta con stakeholder tecnici e di business per definire requisiti funzionali, prototipare rapidamente soluzioni AI e condurre sessioni di user acceptance testing.
\end{itemize}

% ================= NOTABLE PROJECTS =================
\section{Progetti Chiave / Key Projects}
\begin{itemize}
  \item \textbf{Enterprise Multi-Agent Ecosystem via MCP (Model Context Protocol):} Piattaforma open-source di agenti cooperativi orchestrati per eseguire task multi-step su filesystem, database e servizi web esterni, con tracciamento della memoria contestuale e validazione dello stato.
  \item \textbf{Production GraphRAG \& Hybrid Retrieval Engine:} Motore RAG con estrazione automatizzata di nodi/archi su Neo4j combinato con hybrid search su Qdrant, validato tramite pipeline di test continuo RAGAS con punteggi medi di fedeltà superiore al 92\%.
\end{itemize}

% ================= EDUCATION & CERTIFICATIONS =================
\section{Formazione \& Lingue / Education \& Languages}
\begin{itemize}
  \item \textbf{Bachelor of Arts in Philosophy \& Economics} -- Università Ca' Foscari Venezia (2018 -- 2022) \\
  \textit{Focus su logica formale, modelli decisionali, teoria dei giochi ed analisi quantitativa.}
  \item \textbf{Lingue:} Italiano (fluente / operativo), Inglese (C1 professionale, eccellente padronanza scritta e orale).
  \item \textbf{Stato Legale \& Fiscale:} Cittadino extra-UE con Permesso di Soggiorno UE per Soggiornanti di Lungo Periodo a Tempo Indeterminato (diritto al lavoro senza visto in Italia/UE); Partita IVA attiva / regime forfettario pronto.
\end{itemize}

\end{document}
